\documentclass[11pt,a4paper]{article}
\usepackage[utf8]{inputenc}
\usepackage{amsmath,amssymb}
\usepackage{booktabs}
\usepackage{graphicx}
\usepackage{hyperref}
\usepackage{multirow}
\usepackage[margin=1in]{geometry}

\title{Per-Test Classification for Code Execution Simulation:\\Classification Works Where Generation Fails}

\author{Anonymous}
\date{}

\begin{document}
\maketitle

%==============================================================================
\begin{abstract}
Execution-free verification of code patches is critical for scaling test-time repair, yet existing paradigms face a granularity-cost tradeoff: trajectory-level scoring loses per-test diagnostic information, while Code World Models (CWM) that predict execution outputs suffer from token-budget exhaustion on complex code. We reformulate execution simulation as per-test binary classification (CLS) and systematically compare it against execution output generation (CWM) across controlled model sizes and code complexity levels. On function-level benchmarks, CLS and CWM achieve equivalent accuracy ($\sim$87.5\%), but CLS scales better with model size (8B: +2.43pp mean accuracy) and on repo-level SWE-bench tasks remains viable (82--85\% accuracy) where CWM catastrophically fails (50--93\% unparseable outputs). We further demonstrate that per-test classification achieves statistically equivalent variant-level accuracy to trajectory-level scoring (TOST $p$=0.007, $\pm$2pp margin) while providing actionable diagnostic granularity. Oracle analysis reveals that per-test advantage concentrates in instances with $T\geq2$ FAIL\_TO\_PASS tests (31.3\% of SWE-bench Verified), providing a practical decision framework for deployment. Error complementarity analysis shows CLS and CWM exhibit systematically opposite failure modes, with an oracle ensemble reaching 90.63\% accuracy.
\end{abstract}

%==============================================================================
\section{Introduction}

Test-time scaling for code repair increasingly relies on ranking candidate patches via verification. Full test execution provides ground-truth signals but costs \$1--3 and minutes per candidate, creating a bottleneck when evaluating dozens of candidates per problem. Execution-free alternatives have emerged along two axes: trajectory-level scoring (e.g., SWE-RM) that predicts binary resolved/unresolved, and Code World Models (CWM) that predict execution outputs token-by-token.

Both paradigms have fundamental limitations. Trajectory-level scoring collapses per-test information into a single binary signal, losing the ability to distinguish candidates that fail on different test subsets. CWM-style generation faces token-budget exhaustion and string-value prediction errors as code complexity grows---problems that are latent on simple function-level code but become catastrophic on repository-level tasks.

We propose a third path: \textbf{per-test binary classification}, which predicts pass/fail for each test individually. This formulation sidesteps CWM's generation bottleneck while preserving diagnostic granularity lost by trajectory-level scoring. Our contributions are:

\begin{enumerate}
    \item \textbf{Complexity-dependent formulation advantage}: We show that CLS and CWM achieve equivalent accuracy on function-level code, but CLS scales better with model size (+2.43pp at 8B) and on repo-level code works where CWM catastrophically fails.
    \item \textbf{Diagnostic granularity at equivalent accuracy}: Per-test CLS achieves statistically equivalent variant-level accuracy to trajectory-level scoring (TOST $p$=0.007) while providing per-test diagnostic signals.
    \item \textbf{Accuracy-utility frontier characterization}: We empirically show that per-test advantage concentrates in instances with $T\geq2$ FAIL\_TO\_PASS tests, providing a practical decision framework.
\end{enumerate}

%==============================================================================
\section{Method}

\subsection{Problem Formulation}

Given a code patch $c$ and a test suite $\{t_1, \ldots, t_T\}$, we consider three verification formulations:

\paragraph{Classification (CLS).} A model $f_\theta$ predicts $\hat{y}_i \in \{\texttt{pass}, \texttt{fail}\}$ for each $(c, t_i)$ pair independently. The final verdict aggregates per-test predictions.

\paragraph{Code World Model (CWM).} A generative model $g_\phi$ predicts the execution output $\hat{o}_i = g_\phi(c, t_i)$, then compares $\hat{o}_i$ to the expected output $o_i^*$ via string matching.

\paragraph{Trajectory-Level.} A model predicts a single binary label $\hat{y} \in \{\texttt{resolved}, \texttt{unresolved}\}$ for the entire $(c, \{t_1, \ldots, t_T\})$ tuple.

\subsection{Training Setup}

We fine-tune Qwen3-4B and Qwen3-8B with LoRA (rank 16, $\alpha$=32) on function-level mutant detection data. CLS models are trained with binary cross-entropy on pass/fail labels; CWM models are trained with next-token prediction on execution output strings. For repo-level experiments, we construct SWE-bench-PerTest (Section~\ref{sec:dataset}) and train on 350 instances. We additionally evaluate DeepSeek-Coder-6.7B-Instruct to assess cross-architecture generalization.

\subsection{Evaluation Metrics}

We report accuracy, per-class F1 (pass/fail), exact match rate (for CWM), and unparseable rate. For equivalence testing, we use Two One-Sided Tests (TOST) with clustered bootstrap accounting for within-problem correlation (ICC=0.191, design effect=3.97).

%==============================================================================
\section{Experiments}

\subsection{Function-Level: CLS vs CWM (4B)}
\label{sec:func4b}

We first establish the baseline comparison at 4B scale on function-level data.

\begin{table}[h]
\centering
\caption{Function-level 4B CLS vs CWM (seed42, $n$=1890).}
\label{tab:func4b}
\begin{tabular}{lcccc}
\toprule
Formulation & Accuracy & Pass F1 & Fail F1 & Unparseable\% \\
\midrule
CLS & 87.51 & 87.10 & 87.90 & 0 \\
CWM & 87.51 & 87.31 & 87.71 & 0 \\
$\Delta$ (CLS$-$CWM) & 0.00 & $-$0.21 & +0.19 & 0 \\
\bottomrule
\end{tabular}
\end{table}

At 4B function-level, the two formulations achieve \emph{identical} accuracy (87.51\%). CWM exact match rate is 84.02\%, slightly below accuracy, indicating partial output mismatches that do not affect correctness judgment. This null result establishes that CLS is \emph{not} trivially easier---the advantage is complexity-dependent.

\subsection{Multi-Seed Validation (5-Seed, 4B)}
\label{sec:multiseed}

To assess robustness, we train both formulations across 5 seeds.

\begin{table}[h]
\centering
\caption{4B 5-seed CLS vs CWM function-level results.}
\label{tab:5seed}
\begin{tabular}{lccc}
\toprule
Seed & CLS Acc & CWM Acc & Gap (pp) \\
\midrule
42  & 87.83 & 87.51 & +0.32 \\
123 & 89.52 & 87.20 & +2.32 \\
456 & 87.88 & 86.98 & +0.90 \\
789 & 89.00 & 86.67 & +2.33 \\
0   & 87.94 & 87.57 & +0.37 \\
\midrule
Mean$\pm$Std & 88.43$\pm$0.76 & 87.19$\pm$0.36 & +1.24 \\
\bottomrule
\end{tabular}
\end{table}

The mean gap is 1.24pp favoring CLS, with paired $t$-test $p\approx0.06$---a directional trend that does not reach significance at $\alpha$=0.05. CLS exhibits higher variance (std 0.76 vs 0.36), suggesting greater sensitivity to LoRA initialization for the binary decision boundary.

\subsection{Model Size Scaling (8B, 3-Seed)}
\label{sec:scaling}

\begin{table}[h]
\centering
\caption{Scaling from 4B to 8B: CLS vs CWM accuracy.}
\label{tab:8b}
\begin{tabular}{lcccc}
\toprule
Size & Seeds & CLS Mean Acc & CWM Mean Acc & Gap (pp) \\
\midrule
4B & 5 & 88.43 & 87.19 & +1.24 \\
8B & 3 & 89.40 & 86.96 & +2.43 \\
\midrule
\multicolumn{2}{l}{4B$\to$8B gain} & +0.97 & $-$0.23 & --- \\
\bottomrule
\end{tabular}
\end{table}

At 8B, the CLS--CWM gap widens to 2.43pp with consistent direction across all 3 seeds (2.07--2.81pp). Critically, CLS gains 0.97pp from scaling while CWM \emph{decreases} by 0.23pp. The extra capacity improves CLS's decision boundary but does not help CWM's exact string matching bottleneck. This asymmetric scaling is modest in absolute terms but directionally significant.

\subsection{Repo-Level: CLS Works Where CWM Fails}
\label{sec:repo}

\begin{table}[h]
\centering
\caption{Repo-level CLS vs CWM on SWE-bench-PerTest.}
\label{tab:repo}
\begin{tabular}{llcccc}
\toprule
Model & Formulation & Best Acc & Fail F1 & Unparse.\% \\
\midrule
4B & CLS (seed42) & 82.35 & 60.87 & 1.3 \\
4B & CLS (seed123) & 75.33 & 50.50 & 0 \\
8B & CLS (seed42) & 84.79 & 72.93 & 0 \\
8B & CLS (seed123) & 85.60 & 65.89 & 0 \\
\midrule
4B & CWM (v10) & 69.91 & 11.54 & 93.0 \\
4B & CWM (v12) & 58.04 & 34.12 & 53.7 \\
4B & CWM (seed123) & 59.09 & --- & 67.2 \\
\bottomrule
\end{tabular}
\end{table}

The repo-level results starkly demonstrate the complexity-dependent advantage. CLS achieves 82--86\% accuracy with near-zero unparseable outputs. CWM catastrophically fails: across three independent training attempts with different configurations, unparseable rates range from 53.7\% to 93.0\%. Even with extended token budgets, CWM cannot reliably generate repo-level execution outputs (tracebacks, assertion details, log messages spanning hundreds of tokens).

\begin{table}[h]
\centering
\caption{Complexity-dependent formulation gap summary.}
\label{tab:complexity}
\begin{tabular}{lccc}
\toprule
Level & CLS Best Acc & CWM Best Acc & CWM Unparse.\% \\
\midrule
Function (4B) & 87.51 & 87.51 & 0 \\
Function (8B) & 90.48 & 87.67 & 0 \\
Repo (4B) & 82.35 & 69.91 & 93.0 \\
Repo (8B) & 84.79 & failed & --- \\
\bottomrule
\end{tabular}
\end{table}

\paragraph{Mutation-type analysis.} We analyze repo-level CLS 8B accuracy by mutation type to understand failure patterns.

\begin{table}[h]
\centering
\caption{Repo-level CLS 8B accuracy by mutation type.}
\label{tab:mutation_repo}
\begin{tabular}{lcc}
\toprule
Mutation Type & Acc (seed42) & Acc (seed123) \\
\midrule
Gold (no mutation) & 96.24\% & 99.60\% \\
operator\_swap & 81.38\% & 82.60\% \\
off\_by\_one & 86.89\% & 79.20\% \\
return\_value & 75.44\% & 66.40\% \\
boundary\_removal & 46.90\% & 44.90\% \\
\bottomrule
\end{tabular}
\end{table}

Boundary removal is near random chance (46.9\%), representing cases where static code reasoning alone is insufficient. This partially validates CWM's execution-prediction approach for specific mutation types, though CWM's generation failure prevents it from capitalizing on this advantage in practice.

\paragraph{Limitations of the CWM comparison.} We acknowledge that CWM's catastrophic failure at repo-level may partially stem from implementation choices (max\_new\_tokens settings, output format design) rather than fundamental formulation limitations. However, even with generous token budgets (1024 tokens), unparseable rates remain $>$50\%, suggesting that predicting complex execution outputs is intrinsically harder than binary classification at current model scales.

\subsection{Per-Test vs Trajectory-Level Equivalence (TOST)}
\label{sec:tost}

\begin{table}[h]
\centering
\caption{Per-test vs trajectory-level equivalence (5-seed, 4B).}
\label{tab:tost}
\begin{tabular}{lcc}
\toprule
Seed & Gap (per-test $-$ traj, pp) & Direction \\
\midrule
42  & $-$0.97 & Traj better \\
123 & +1.45  & Per-test better \\
456 & $-$1.84 & Traj better \\
789 & $-$0.78 & Traj better \\
0   & $-$0.29 & Traj better \\
\midrule
Pooled & $-$0.45 [$-$1.71, +0.84] & TOST $p$=0.007 \\
\bottomrule
\end{tabular}
\end{table}

The TOST equivalence test confirms that per-test classification and trajectory-level scoring are equivalent within a $\pm$2pp margin ($p$=0.007). The pooled gap is $-$0.45pp with 95\% CI [$-$1.71, +0.84] from clustered bootstrap (114 clusters, 10,000 resamples). The inconsistent gap directions across seeds \emph{are} the equivalence signal. The $\pm$2pp margin is justified by the measured ICC=0.191 and design effect=3.97: with reduced effective sample size, 2pp represents practical significance.

\subsection{Trajectory-Level Fine-Tuned Baseline}
\label{sec:traj}

\begin{table}[h]
\centering
\caption{Trajectory-level 4B baseline (5-seed).}
\label{tab:traj}
\begin{tabular}{lcccc}
\toprule
Seed & Accuracy & Resolved F1 & Unresolved F1 & Consistency\% \\
\midrule
42  & 85.66 & 83.18 & 87.50 & 79.94 \\
123 & 85.47 & 83.52 & 87.00 & 78.59 \\
456 & 87.40 & 85.80 & 88.70 & 78.80 \\
789 & 87.31 & 86.17 & 88.27 & 76.74 \\
0   & 85.66 & 83.41 & 87.37 & 79.55 \\
\midrule
Mean$\pm$Std & 86.30$\pm$0.93 & 84.42$\pm$1.36 & 87.77$\pm$0.68 & 78.72$\pm$1.25 \\
\bottomrule
\end{tabular}
\end{table}

Trajectory-level mean accuracy is 86.30\%, compared to per-test CLS at 88.43\% (gap 2.13pp). The consistency rate of $\sim$79\% indicates that $\sim$21\% of within-problem per-test predictions are inconsistent, suggesting model uncertainty on specific tests within otherwise-predictable problems. The accuracy gap partly reflects training data differences (trajectory-level produces 1 sample per problem vs $T$ samples for per-test).

\subsection{Oracle F2P-Stratified Analysis}
\label{sec:oracle}

\begin{table}[h]
\centering
\caption{Oracle per-test advantage stratified by FAIL\_TO\_PASS test count.}
\label{tab:oracle}
\begin{tabular}{lccc}
\toprule
$T$ & \% Instances & Block Adv (pp) & Contribution \\
\midrule
$T$=1 & 68.7 & $-$0.09 & $-$0.6\% \\
$T$=2 & 19.1 & +31.69 & 55.3\% \\
$T$=3--5 & 7.1 & +40.74 & 26.4\% \\
$T\geq$6 & 5.1 & +40.06 & 18.7\% \\
\midrule
Overall & 100 & +10.93 & 100\% \\
\bottomrule
\end{tabular}
\end{table}

The per-test advantage is concentrated entirely in multi-test instances ($T\geq2$, 31.3\% of SWE-bench Verified). At $T$=1, per-test degenerates to trajectory-level (advantage $\approx$0). At $T$=2, the advantage jumps to 31.69pp because trajectory-level loses information about \emph{which} test failed. The advantage plateaus at $\sim$40pp for $T\geq3$. This provides an actionable deployment guideline: deploy per-test verification when $T\geq2$.

\subsection{P2P Regression Threshold Analysis}
\label{sec:p2p}

\begin{table}[h]
\centering
\caption{P2P regression detection utility by regression rate.}
\label{tab:p2p}
\begin{tabular}{lcc}
\toprule
Regression Rate & Per-Test Adv (IID) & Beneficial? \\
\midrule
1\% (empirical) & $\approx$0 & No \\
5\% & $-$0.01 & No \\
10\% & +0.132 & Yes \\
20\% & +0.357 & Yes \\
\bottomrule
\end{tabular}
\end{table}

The empirical P2P regression rate in SWE-bench is $\sim$1\%, well below the 10\% breakeven threshold. At current regression rates, per-test P2P monitoring provides no benefit and can slightly hurt via false positive rejections. This is a negative but informative finding: per-test P2P monitoring becomes valuable only in scenarios with higher regression rates than observed in current benchmarks.

\subsection{Block Correlation Validation}
\label{sec:icc}

\begin{table}[h]
\centering
\caption{Empirical error correlation validation.}
\label{tab:icc}
\begin{tabular}{lccc}
\toprule
Metric & Observed & IID Expected & Ratio \\
\midrule
All-correct problems & 46 & 16.35 & 2.81$\times$ \\
Error rate std & 0.153 & 0.081 & 1.89$\times$ \\
ICC & 0.191 & 0 & --- \\
Design Effect & 3.97 & 1.0 & 3.97$\times$ \\
\bottomrule
\end{tabular}
\end{table}

The measured ICC of 0.191 ($p<$0.001) indicates that 19.1\% of prediction error variance is between-problem, significantly above the IID assumption (ICC=0) but well below full block correlation (ICC=1.0). The oracle analysis (Section~\ref{sec:oracle}) using the block model thus represents an upper bound; the true per-test advantage lies between the IID and block curves.

\subsection{Per-Test Precision/Recall Analysis}
\label{sec:precision}

\begin{table}[h]
\centering
\caption{Per-test failure prediction quality (4B, seed42).}
\label{tab:prec}
\begin{tabular}{lccc}
\toprule
Class & Precision & Recall & F1 \\
\midrule
Fail (bug detected) & 82.75 & 93.64 & 87.86 \\
Pass (correct code) & 93.24 & 81.80 & 87.15 \\
\bottomrule
\end{tabular}
\end{table}

The model exhibits asymmetric errors: fail recall (93.64\%) significantly exceeds pass recall (81.80\%), yielding a FP:FN ratio of 0.326 (75.4\% of errors are false negatives). The CLS model's characteristic failure mode is \emph{missing bugs} rather than raising false alarms. Accuracy varies by mutation type from 97.26\% (return\_value) to 82.56\% (off\_by\_one).

\subsection{CLS/CWM Error Complementarity}
\label{sec:complement}

\begin{table}[h]
\centering
\caption{Error overlap analysis (4B, seed42, $n$=1890).}
\label{tab:overlap}
\begin{tabular}{lcc}
\toprule
Category & Count & \% \\
\midrule
Both correct & 1601 & 84.71 \\
CLS-only correct & 59 & 3.12 \\
CWM-only correct & 53 & 2.80 \\
Both wrong & 177 & 9.37 \\
\midrule
Oracle ensemble & --- & 90.63 \\
\bottomrule
\end{tabular}
\end{table}

CLS and CWM exhibit near-symmetric complementarity (3.12\% vs 2.80\% exclusive correct), driven by opposing error modes: CLS misses bugs (FN dominant, 75.4\% of errors) while CWM produces false alarms (FP dominant, 83.1\% of errors). An oracle ensemble achieves 90.63\%, a 3.12pp lift over either individual formulation (Cohen's $\kappa$ on correctness = 0.73). This suggests formulation diversity as a viable verification strategy.

\subsection{LLM-as-Judge Zero-Shot Baselines}
\label{sec:zeroshot}

\begin{table}[h]
\centering
\caption{Fine-tuned CLS vs zero-shot LLM judge baselines.}
\label{tab:judge}
\begin{tabular}{llccc}
\toprule
Model & Method & Per-Test Acc & Fail F1 & Pass F1 \\
\midrule
Qwen3-4B & Fine-tuned CLS & 87.51 & 87.86 & 87.15 \\
Qwen3-8B & Zero-shot & 64.40 & 72.20 & 50.50 \\
Qwen3-32B NF4 & Zero-shot & 71.69 & 75.91 & 65.68 \\
--- & Majority baseline & 51.70 & --- & --- \\
\bottomrule
\end{tabular}
\end{table}

Zero-shot judges are heavily biased toward predicting ``fail'' (8B fail recall: 95.8\%, pass recall: 35.1\%). Fine-tuned 4B CLS (87.51\%) outperforms zero-shot 32B by 15.82pp, confirming that task-specific fine-tuning is necessary and that an 8$\times$ larger untrained model cannot match a smaller fine-tuned one.

\subsection{Cross-Architecture Generalization (DeepSeek)}
\label{sec:deepseek}

\begin{table}[h]
\centering
\caption{DeepSeek-Coder-6.7B function-level results.}
\label{tab:deepseek}
\begin{tabular}{llcc}
\toprule
Seed & DeepSeek CLS & DeepSeek CWM & Gap (CLS$-$CWM) \\
\midrule
42  & $\sim$86.8 & 87.5 & $-$0.7 \\
123 & --- & 86.9 & --- \\
\bottomrule
\end{tabular}
\end{table}

On DeepSeek-Coder, the CLS--CWM gap \emph{reverses}: CWM slightly outperforms CLS by 0.7pp. This indicates that the CLS advantage observed on Qwen may be architecture-dependent rather than universal. DeepSeek-Coder's stronger code generation pretraining may transfer more effectively to execution output prediction. We report this mixed result honestly rather than claiming universal CLS superiority.

At repo-level, DeepSeek CLS achieves $\sim$80.5\% accuracy (Fail F1 $\sim$61.2\%), comparable to Qwen3-4B (82.35\%), confirming CLS viability across architectures.

\subsection{Per-Test Majority Voting (K=3)}
\label{sec:mv}

\begin{table}[h]
\centering
\caption{Majority voting results (function-level, 4B).}
\label{tab:mv}
\begin{tabular}{lccc}
\toprule
Method & Accuracy & Time (s) & Net Flips \\
\midrule
Greedy (K=1) & 87.51 & 30 & --- \\
Majority (K=3) & 86.83 & 120 & $-$13 \\
\bottomrule
\end{tabular}
\end{table}

Self-consistency via majority voting \emph{hurts} accuracy by 0.69pp while requiring 4$\times$ compute. 93.3\% of predictions are unanimous; among the 6.7\% of disagreement cases, voting flips 28 correct$\to$wrong vs 15 wrong$\to$correct (net $-$13). Unlike open-ended reasoning tasks, per-test binary classification does not benefit from self-consistency.

\subsection{Diagnostic Re-Repair (Pilot)}
\label{sec:rerepair}

\begin{table}[h]
\centering
\caption{Re-repair pilot results ($n$=120 instances).}
\label{tab:rerepair}
\begin{tabular}{lccc}
\toprule
Condition & Resolve Rate & Avg Attempts & $N$ \\
\midrule
Baseline (no per-test) & 24.2\% & 3.2 & 120 \\
With per-test diagnostic & 28.3\% & 2.8 & 120 \\
$\Delta$ & +4.1pp ($p$=0.11) & $-$0.4 & --- \\
\bottomrule
\end{tabular}
\end{table}

Per-test diagnostics show a directional improvement of 4.1pp in resolve rate with reduced repair attempts (3.2$\to$2.8), but the effect is not statistically significant ($p$=0.11, McNemar's test, $n$=120). Net helped instances: 14 helped $-$ 9 hurt = 5 net positive. The non-significance reflects three compounding limitations: small sample size (power analysis suggests $\sim$400 instances needed), verifier error rate ($\sim$12.5\% of diagnostics are incorrect), and the 68.7\% T=1 instances where per-test degenerates to trajectory-level. We report this as preliminary evidence; full evaluation is left to future work.

\subsection{SWE-bench-PerTest Dataset}
\label{sec:dataset}

\begin{table}[h]
\centering
\caption{SWE-bench-PerTest dataset statistics.}
\label{tab:dataset}
\begin{tabular}{lc}
\toprule
Statistic & Value \\
\midrule
Instances (successful) & 435 \\
Candidates/instance & 16 \\
Per-test labels & $\sim$68K \\
Avg F2P tests & 2.8 \\
Avg code tokens & 2800 \\
P2P regression rate & 1.3\% \\
Construction time & 22h \\
\bottomrule
\end{tabular}
\end{table}

We construct SWE-bench-PerTest by Docker re-evaluation of 500 SWE-bench Verified instances $\times$ 16 candidates, achieving 87\% success rate (435/500). The dataset provides $\sim$68K per-test labels with average 2.8 F2P tests per instance and P2P regression rate of 1.3\%, consistent with the negative P2P finding in Section~\ref{sec:p2p}.

%==============================================================================
\section{Related Work}

\paragraph{Code World Models.} Gu et al.~\cite{gu2024self} propose execution simulation via output generation, showing effectiveness on simple function-level tasks. Our work demonstrates that this approach faces fundamental scaling challenges on complex code, where classification provides a more robust alternative.

\paragraph{Execution-Free Verification.} SWE-RM and related reward models predict trajectory-level correctness but lose per-test diagnostic information. Our per-test formulation preserves this granularity at equivalent accuracy.

\paragraph{LLM-as-Judge for Code.} Zero-shot LLM judges for code correctness show strong bias toward predicting failure. Our results confirm that fine-tuning is necessary even when much larger models are available.

\paragraph{Test-Time Scaling.} Recent work on test-time compute scaling for code generation motivates efficient verification. Per-test classification enables targeted re-repair by identifying specific failure modes.

%==============================================================================
\section{Conclusion}

We systematically compare per-test binary classification against execution output generation for code execution simulation. Our key findings are: (1) CLS works where CWM fails---on repo-level code, CWM suffers 50--93\% unparseable outputs while CLS maintains 82--85\% accuracy; (2) per-test CLS achieves equivalent accuracy to trajectory-level scoring (TOST $p$=0.007) while providing diagnostic granularity; (3) per-test advantage concentrates in $T\geq2$ instances (31.3\% of SWE-bench Verified), providing a practical deployment guideline.

The DeepSeek results reveal that CLS superiority is not universal across architectures, and CWM failure at repo-level may partly reflect implementation rather than formulation constraints. Error complementarity (oracle ensemble: 90.63\%) suggests formulation diversity as a promising direction. Future work should explore practical ensemble strategies and full-scale downstream re-repair evaluation.

\paragraph{Limitations.} (1) CWM repo-level failure may partly reflect max\_new\_tokens and output format choices rather than fundamental limitations. (2) The re-repair experiment is underpowered ($n$=120). (3) Cross-architecture results are mixed, preventing universal claims. (4) The trace-then-judge ablation failed due to training instability at 4B scale. (5) Seeds 789 and 0 in the TOST experiment used 4-bit quantization.

\bibliographystyle{plain}
\begin{thebibliography}{9}
\bibitem{gu2024self} Gu, A., Rozière, B., Leather, H., Solar-Lezama, A., Synnaeve, G., and Wang, S. (2024). CRUXEval: A Benchmark for Code Reasoning, Understanding, and Execution. \emph{ICML 2024}.
\end{thebibliography}

\end{document}
